\scanpage{224}
\mspage{79}

Les parties closes sont les sous-groupes fermés.

3) Pour deux espaces vectoriels ``\unsure{en} dualité'' sur un corps valué complet, tel $\mathbb R$ ou $\mathbb C$, la relation

\[
 |\langle x,x'\rangle|\leq 1.
\]

Dans le cas réel ou complexe, les parties closes sont les parties convexes \unsure{cerclées} faiblement fermées

4) Pour les espaces vectoriels réels, la relation de semi-polarité

\[
 \langle x,x'\rangle\deleted{\unsure{>|\geq}}\unsure{\geq|\geqq}-1.
\]

Les parties closes sont les parties convexes \unsure{fblt.} fermées contenant l'origine.

5) Soit $K$ un corps, \inserted{(pour simplifier)} $L$ une extension galoisienne de groupe $G$, Relation entre $L$ et $G$

\[
 \sigma.x=x\qquad\text{i.e. }x\in L^\sigma.
\]

Les parties closes de $L$ sont les sous-extensions, les parties closes de $G$ les sous-groupes de $G$

Il y a des \deleted{\unsure{des|des deg.}} analogues divers \unsure{dans des|dans les|de ces} cas \unsure{inséparables} (à \unsure{coup\deleted{t}s|coups|compte} de $p$-algèbres de Lie par exemple), et des énoncés de descente en tous genres.
% archive sheet number: 224
